% !TeX root = ../../main.tex
\subsection{感嘆文の種類}
    感嘆文とは、驚き、喜び、悲しみ、怒りなどの\uwave{感情や感動をストレートに表現する文}のことである。
    感嘆文の大きな特徴の1つとして、文末に\fitblankbf[]{!}がつくことが挙げられる。

    感嘆文は\fitblankbf[]{what} を使った感嘆文と\fitblankbf[]{how} を使った感嘆文の2パターンしかない。

\subsection{whatを使った感嘆文}

\begin{theorembox}{\textbf{whatを使った感嘆文}}
    \fitblankbfresize{%
      \begin{tabular}{@{}c@{\hspace{0.15em}+\hspace{0.15em}}c@{\hspace{0.15em}+\hspace{0.15em}}c@{\hspace{0.15em}+\hspace{0.15em}}c@{\hspace{0.15em}+\hspace{0.15em}}c@{\hspace{0.15em}+\hspace{0.15em}}c@{\hspace{0.15em}+\hspace{0.15em}}c@{}}%
        \fitblankbf[]{What} & \fitblankbfshort{a[an]} & \fitblankbf[]{形容詞} & \fitblankbf[]{名詞} & \fitblankbf[]{主語} & \fitblankbf[]{動詞} & \fitblankbfshort{!}%
      \end{tabular}%
    }
\vspace{1.5em}

「〜はなんて \fitblankbf[]{形容詞}な ・・・ でしょう。」

\end{theorembox}
\vspace{0.5em}

whatを使った感嘆文は\fitblankbf[]{7つ}のブロックから構成されていることを意識しよう!

\begin{enumerate}
    \item What a cool math teacher he is!
    \dialogueblank{彼はなんてカッコいい数学の先生でしょう。}
    \vspace{0.5em}
    \item これはなんて大きなケーキでしょう。
    \dialogueblank{What a big cake this is!}
    \vspace{1em}
    \item \fitblankbf[]{彼はなんてカッコいい数学の先生でしょう。}
    \dialogueblank{What a cool math teacher he is!}
\end{enumerate}

\subsection{howを使った感嘆文}

    \begin{theorembox}{\textbf{howを使った感嘆文}}
        \fitblankbfresize{%
          \begin{tabular}{@{}c@{\hspace{0.15em}+\hspace{0.15em}}c@{\hspace{0.15em}+\hspace{0.15em}}c@{\hspace{0.15em}+\hspace{0.15em}}c@{\hspace{0.15em}+\hspace{0.15em}}c@{}}%
            \fitblankbf[]{How} & \fitblankbf[]{形容詞 [ 副詞 ]} & \fitblankbf[]{主語} & \fitblankbf[]{動詞} & \fitblankbfshort{!}%
          \end{tabular}%
        }
    \vspace{1.5em}

        「〜はなんて \fitblankbf[]{形容詞 [ 副詞 ]} なのでしょう。」

    \end{theorembox}
    \vspace{1em}

    howを使った感嘆文は\fitblankbf[]{5つ}のブロックから構成されていることを意識しよう!

    \begin{enumerate}
        \item How cute !
        \dialogueblank{なんてかわいいのでしょう。}
    \vspace{0.5em}
        \item この家はなんて古いのでしょう。
        \dialogueblank{How old this house is!}
    \vspace{1em}
        \item \fitblankbf[]{なんてかわいいのでしょう。}
        \dialogueblank{How cute !}
    \end{enumerate}
